\documentclass[12pt]{article}
\usepackage{graphicx} % Required for inserting images
\usepackage[margin=0.5in]{geometry}
\usepackage{amsmath, amssymb}
\usepackage{mathrsfs}


\usepackage{soul}


% Dr. David Brown Commands
\newcommand{\ds}{\displaystyle}
\newcommand{\R}{\mathbb{R}}
\newcommand{\N}{\mathbb{N}}
\newcommand{\Q}{\mathbb{Q}}
\newcommand{\C}{\mathbb{C}}


% Commands to get script r
\usepackage{calligra}
\DeclareMathAlphabet{\mathcalligra}{T1}{calligra}{m}{n}
\DeclareFontShape{T1}{calligra}{m}{n}{<->s*[2.2]callig15}{}
\newcommand{\scripty}[1]{\ensuremath{\mathcalligra{#1}}}

% Unit commands
\newcommand{\degree}{\text{°}}
\newcommand{\unitSpeed}{\frac{\text{m}}{\text{s}}}
\newcommand{\unitAcceleration}{\frac{\text{m}}{\text{s}^2}}

% Constant commands
\newcommand{\avogadro}{6.022\times10^{23}}

\newcommand{\boltzmannConstK}{1.381\times10^{-23}\frac{\text{J}}{\text{K}}}
\newcommand{\boltzmannConstInEV}{8.617\times10^{-5}\frac{\text{eV}}{\text{K}}}
\newcommand{\planckConst}{6.626\times10^{-34}\text{J}\cdot\text{s}}

\newcommand{\atmP}{1.01325\times10^5\,\text{Pa}}

% Commands to easily write partial derivatives
\newcommand{\partDeriv}[2]{\frac{\partial #1}{\partial #2}}
\newcommand{\doublePartDeriv}[2]{\frac{\partial^2 #1}{\partial^2 #2}}


\title{Problem 7.13}
\author{Gabriel Decker}


\begin{document}



\maketitle
\begin{flushleft}


For this problem, we will determine the average occupancy of a state for a system of bosons, varying the energy levels.
To do this, we'll use the bose-einstein distribution.
\begin{equation}
    \overline{n}_\text{BE}=\frac{1}{e^{(\epsilon-\mu)/kT}-1}
\end{equation}
where $\epsilon$ is the energy of the state.\\~\\

We will also find the probability of finding $n=1,2,3$ bosons in a state.
Recall that the probability of a number of particles being in a state is the following.
\begin{equation}
    \mathcal{P}(n)=\frac{1}{\mathcal{Z}}e^{-n(\epsilon-\mu)/kT}
\end{equation}
where for bosons, $\mathcal{Z}$ is the following.
\begin{equation}
    \mathcal{Z}=\frac{1}{1-e^{-(\epsilon-\mu)/kT}}
\end{equation}
We can combine these equations into the following.
$$\mathcal{P}(n)=\frac{1}{\mathcal{Z}}e^{-n(\epsilon-\mu)/kT}$$
$$\implies\mathcal{P}(n)=(1-e^{-(\epsilon-\mu)/kT})e^{-n(\epsilon-\mu)/kT}$$\\~\\

We can substitute $x=(\epsilon-\mu)/kT$ to get the following.
$$\implies\mathcal{P}(n)=(1-e^{-x})e^{-nx}$$
and
$$\overline{n}_\text{BE}=\frac{1}{e^{x}-1}$$\\~\\

We'll find the occupancies and probabilities for the following values of $\epsilon$.
\begin{enumerate}
    \item[a.] $\epsilon=\mu+0.001\,\text{eV}$:\\
    \item[b.] $\epsilon=\mu+0.01\,\text{eV}$:\\
    \item[c.] $\epsilon=\mu+0.1\,\text{eV}$:\\
    \item[d.] $\epsilon=\mu+1\,\text{eV}$:\\
\end{enumerate}

We'll be doing calculations at $T=300\,\text{K}$.
We can further substitue $x$ by representing the offsets from $\mu$ as $\delta$.
$$x=(\epsilon-\mu)/kT$$
$$\implies x=(\mu+\delta-\mu)/kT$$
$$\implies x=\delta/kT$$
We can go ahead and plug in $k$ and $T$ since that is always the same.
$$\implies x=\delta/0.025851\,\text{eV}$$\\~\\


\textbf{Part a}\\
For this part, $\delta=0.001\,\text{eV}$.
Plugging this into the above equations, we get the following.
$$\overline{n}_\text{BE}=25.354$$
$$\mathcal{P}(0)=0.0379$$
$$\mathcal{P}(1)=0.0365$$
$$\mathcal{P}(2)=0.0351$$
$$\mathcal{P}(3)=0.0338$$

\textbf{Part b}\\
For this part, $\delta=0.01\,\text{eV}$.
Plugging this into the above equations, we get the following.
$$\overline{n}_\text{BE}=2.117$$
$$\mathcal{P}(0)=0.321$$
$$\mathcal{P}(1)=0.218$$
$$\mathcal{P}(2)=0.148$$
$$\mathcal{P}(3)=0.101$$

\textbf{Part c}\\
For this part, $\delta=0.1\,\text{eV}$.
Plugging this into the above equations, we get the following.
$$\overline{n}_\text{BE}=0.021$$
$$\mathcal{P}(0)=0.979$$
$$\mathcal{P}(1)=0.0205$$
$$\mathcal{P}(2)=0.000427$$
$$\mathcal{P}(3)=0.000009$$

\textbf{Part d}\\
For this part, $\delta=1\,\text{eV}$.
Plugging this into the above equations, we get the following.
$$\overline{n}_\text{BE}=1.58\times10^{-17}$$
$$\mathcal{P}(0)=1$$
$$\mathcal{P}(1)=1.58\times10^{-17}$$
$$\mathcal{P}(2)=2.51\times10^{-34}$$
$$\mathcal{P}(3)=3.98\times10^{-51}$$





\end{flushleft}

\end{document}
